\documentclass[11pt]{article}
\usepackage[margin=25mm]{geometry}
\usepackage{fontspec,unicode-math}
\setmainfont{Latin Modern Roman}
\setmathfont{Latin Modern Math}
\usepackage{amsmath,amssymb,mathtools}
\usepackage{xurl,hyperref}
\hypersetup{colorlinks=true,urlcolor=blue,linkcolor=blue}
\setcounter{secnumdepth}{0}
\setlength{\parindent}{0pt}
\setlength{\parskip}{5pt}
\setlength{\emergencystretch}{3em}
\newcommand{\sref}[2]{\hyperlink{source:#1:#2}{[S#2]}}
\newcommand{\cataloguescope}[1]{\paragraph{Source question.}#1}
\newcommand{\floor}[1]{\left\lfloor#1\right\rfloor}
\begin{document}
\section{Global Observability from Negative Gramian Curvature}\label{statement}

\subsection{Definitions and mathematical statement}

\paragraph{Shared notation.}Standard number systems are denoted by $\mathbb N,\mathbb Z,\mathbb Q,\mathbb R,\mathbb C$; index ranges, computational models and specialized notation follow the cited problem source and the definitions below.


$\Phi(t)$ is the linearized state-transition matrix and $H(t)$ the output derivative along the trajectory. A uniformly positive observability Gramian defines a Riemannian metric; global observability means distinct initial states give distinct full output trajectories.

In the LOCAL model vertices communicate synchronously along graph edges; messages and local computation are unrestricted and complexity is the number of rounds. After $t$ rounds a vertex can depend only on its radius-$t$ neighborhood and any supplied identifiers or randomness.

A state-space system evolves as $\dot x=f(x,u)$ with output $y=h(x,u)$. Model reduction seeks a lower-dimensional system preserving specified input-output properties. Observability distinguishes initial states from their output trajectories, while controllability concerns reachable states under allowed inputs.
\paragraph{Statement recorded in the supplied source.}
For the observed system $\dot{x}=f(x)$, $y=h(x)$ on an $n$-dimensional manifold, let $P(x_0)=\int_0^T\Phi(t)^TH(t)^TH(t)\Phi(t)\,dt$ be the local observability Gramian along the trajectory from $x_0$. If $P(x_0)$ is uniformly positive definite and bounded, does uniform negative curvature of the resulting Riemannian metric, together with geodesic completeness, imply global observability on $[0,T]$? \sref{statement}{2}
\paragraph{Source status.}
The supplied catalogue labels this record open. Its source-status metadata was checked against the supplied snapshot on 2026-10-11. This import does not independently establish current open status or mathematical correctness. The catalogue literature assessment is dated 2026-08-21.
\subsection{Short English statement}
Determine whether the stated completeness and negative-curvature assumptions promote local observability to global observability.
\subsection{Sources}
\begin{itemize}
\item[\hypertarget{source:statement:1}{S1}] Ulam AI and the UnsolvedMath Contributors, UnsolvedMath: A Curated Collection of Open Mathematics Problems, catalogue entry. Catalogue curation and metadata: CC BY 4.0. The supplied source question is retained; definitions and the English summary are editorial additions.
\url{https://huggingface.co/datasets/ulamai/UnsolvedMath}.
\item[\hypertarget{source:statement:2}{S2}] Control Theory: On the Way to New Application Fields (2009), Oberwolfach Reports, DOI 10.4171/owr/2009/11, pp. 593--680
\url{https://doi.org/10.4171/owr/2009/11}.
\textit{Underlying Oberwolfach Reports material retains its source-specific licence, including CC BY-SA 4.0 where applicable; see THIRD\_PARTY\_NOTICES.md.}
\end{itemize}
\paragraph{Preparation and review.}
Original question attributed to the cited source and UnsolvedMath contributors. Codex (GPT-6; exact serving revision not exposed) assisted with definitions, English summaries, source inspection and formatting. No new human mathematical authorship or independent proof verification is asserted. Catalogue identity checked against the complete supplied registry and existing repository statements; known flagged repetitions and obvious identical conjectures excluded. Full mathematical-equivalence/literature review has not been performed. Statement preparation and final review followed contribution prompts C and D; automated checks verify structure and identity, not open status.
\label{end:statement}
\end{document}
